The following prompt was used identically for both Claude 3.5 Sonnet (Anthropic) and GPT-4o (OpenAI) classifications. Each model received the prompt once per narrative, with temperature set to $0.0$ (deterministic output) and a maximum output length of 300 tokens. No product type or company name was passed to the classifier, preventing hypothesis leakage. The prompt consists of a system prompt defining the construct framework and a user prompt specifying the required output format.

\subsection{System Prompt}

\begin{quote}
\begin{small}
\begin{verbatim}
You are a research assistant classifying consumer financial
complaints for an academic study on trust in financial
relationships.

The study distinguishes two modes of trust failure:

EVALUATIVE TRUST FAILURE: The consumer's complaint is primarily
about unmet performance expectations, functional problems, or
verification failures. The consumer treats the financial
institution as a service provider that failed to deliver. Language
focuses on: incorrect charges, broken processes, unmet promises,
poor service quality, failure to perform as advertised. The
emotional register is frustration, dissatisfaction, or anger at
incompetence. The consumer's implicit stance is: "You didn't do
what you were supposed to do."

DELEGATIVE TRUST FAILURE: The consumer's complaint reflects
betrayal of a relationship in which they had granted the
institution discretionary authority or relied on it as a fiduciary.
Language focuses on: betrayal, abandonment, broken loyalty, abuse
of the relationship, violation of good faith, taking advantage of
vulnerability. The emotional register is moral outrage, a sense of
personal violation, or grief at a relationship destroyed. The
consumer's implicit stance is: "I trusted you and you betrayed me."

UNCLASSIFIABLE: The complaint does not clearly reflect either trust
mode. This includes: template/form letters, pure factual disputes
with no emotional content, identity theft reports where the
consumer has no relationship with the institution, complaints too
brief or unclear to assess, or complaints that mix both modes
without a dominant one.

Important: Many complaints will be unclassifiable. That is
expected. Do not force a classification. Only assign EVALUATIVE or
DELEGATIVE when the language clearly reflects one mode.
\end{verbatim}
\end{small}
\end{quote}

\subsection{User Prompt Template}

For each narrative, the following user prompt was appended after the system prompt, with the complaint text inserted in place of the placeholder:

\begin{quote}
\begin{small}
\begin{verbatim}
Classify this consumer financial complaint. Respond in exactly
this JSON format:

{
  "classification": "EVALUATIVE" | "DELEGATIVE" | "UNCLASSIFIABLE",
  "confidence": <float 0.0 to 1.0>,
  "primary_evidence": "<one sentence quoting or paraphrasing the
    specific language that drove your classification>",
  "secondary_signals": "<one sentence noting any additional
    signals, or 'none'>",
  "delegative_score": <float 0.0 to 1.0>,
  "evaluative_score": <float 0.0 to 1.0>
}

The delegative_score and evaluative_score are independent
dimensions, not a zero-sum scale. A complaint can score high on
both (mixed), low on both (unclassifiable), or high on one and
low on the other (clean classification).

COMPLAINT:
"""
[Narrative text inserted here]
"""
\end{verbatim}
\end{small}
\end{quote}

\subsection{Classification Parameters}

Both models were configured with identical parameters to ensure comparability. \Cref{tab:classify_params} summarizes the classification pipeline settings.

\begin{table}[ht]
\centering
\caption{Classification pipeline parameters.}
\label{tab:classify_params}
\begin{tabular}{lll}
\toprule
Parameter & Claude 3.5 Sonnet & GPT-4o \\
\midrule
Model version & \texttt{claude-3-5-sonnet-20241022} & \texttt{gpt-4o} \\
Temperature & 0.0 & 0.0 \\
Max output tokens & 300 & 300 \\
Concurrent requests & 20 & 5 \\
Retry logic & Exponential backoff (5 attempts) & Exponential backoff (5 attempts) \\
Narratives classified & 8,943 & 8,943 \\
\bottomrule
\end{tabular}
\end{table}
